\chapter*{W}
\label{chap:w}

\term{Waelbroeck Space}
A Waelbroeck space is a topological vector space designed to support a useful theory of bounded operators, often without local convexity. \par\noindent\textbf{Applications:} such spaces occur in generalized-function and distribution theory.

\term{Walk (Graph Theory)}
A walk in a graph is an alternating sequence of vertices and incident edges. Unlike a path, it may repeat vertices or edges. \par\noindent\textbf{Applications:} walks model routes, random processes, search procedures, and network communication.

\term{Wallis Product}
The Wallis product is the convergent infinite product $\frac21\cdot\frac23\cdot\frac43\cdot\frac45\cdots=\frac\pi2$. \par\noindent\textbf{Applications:} it provides a classical product for $\pi$ and illustrates convergence of infinite products.

\term{Walsh--Hadamard Transform}
The Walsh--Hadamard transform expands data in piecewise-constant Walsh functions using only additions and subtractions. \par\noindent\textbf{Applications:} it supports fast signal processing, digital communication, coding, and quantum algorithms.

\term{W-algebra}
A W-algebra is an algebra of operators encoding generalized symmetries, often in conformal field theory. \par\noindent\textbf{Applications:} W-algebras support the study of higher-spin fields, integrable systems, and representation theory.

\term{Ward Identity}
A Ward identity is a relation among correlation functions that follows from a symmetry of a quantum theory. \par\noindent\textbf{Applications:} Ward identities constrain quantum calculations and express conservation laws related to Noether's theorem.

\term{Waring's Problem}
Waring's problem asks whether every positive integer is a sum of a bounded number of $k$th powers for each fixed $k$. \par\noindent\textbf{Applications:} it is a central problem in additive number theory and the study of representations of integers.

\term{Wave Equation}
The wave equation models propagation, commonly as $u_{tt}=c^2\nabla^2u$. \par\noindent\textbf{Applications:} it describes sound, vibrations, elastic waves, electromagnetic waves, and acoustics.

\term{Wavelength}
Wavelength is the distance between successive points with the same phase, such as neighboring crests. \par\noindent\textbf{Applications:} it determines wave frequency, color, resolution, and antenna behavior.

\term{Wavelet Transform}
A wavelet transform decomposes a signal into shifted and scaled copies of a localized wavelet. It records both time and frequency behavior. \par\noindent\textbf{Applications:} wavelets support image compression, denoising, anomaly detection, and signal analysis.

\term{Weak Convergence}
Weak convergence means that every continuous linear functional has the same limit on the sequence as on the candidate limit. \par\noindent\textbf{Applications:} it provides compactness tools for infinite-dimensional optimization and PDEs.

\term{Weak-* Convergence}
Weak-* convergence in a dual space is tested against vectors of the original space rather than all elements of the double dual. \par\noindent\textbf{Applications:} it supports compactness results such as Banach--Alaoglu and convergence of measures.

\term{Weak Derivative}
A weak derivative is defined through integration by parts against smooth test functions, even when the ordinary derivative does not exist. \par\noindent\textbf{Applications:} weak derivatives define Sobolev spaces and allow nonsmooth PDE solutions.

\term{Weak Solution}
A weak solution satisfies a differential equation after integration against suitable test functions, so it may have fewer classical derivatives. \par\noindent\textbf{Applications:} weak solutions are fundamental in PDE theory, finite elements, and fluid mechanics.

\term{Weak Topology}
A weak topology is the coarsest topology making the chosen continuous linear functionals continuous. Its convergence is tested by those functionals. \par\noindent\textbf{Applications:} weak topologies provide compactness and convergence methods in functional analysis.

\term{Weakly Connected Component}
A weakly connected component of a directed graph is a maximal subgraph that becomes connected when edge directions are ignored. \par\noindent\textbf{Applications:} it identifies network regions and communication components.

\term{Weighted Graph}
A weighted graph assigns a numerical weight to each edge, such as a cost, distance, or capacity. \par\noindent\textbf{Applications:} weighted graphs model transport, communication, scheduling, and logistics.

\term{Weighted Mean}
A weighted mean is an average in which observations contribute according to weights: $\sum_iw_ix_i/\sum_iw_i$. \par\noindent\textbf{Applications:} weighted means combine measurements with different reliability and support survey analysis.

\term{Weil Conjecture}
The Weil conjectures are four theorems describing point counts and zeta functions of varieties over finite fields. \par\noindent\textbf{Applications:} they connect algebraic geometry, number theory, and cohomology.

\term{Weil Pairing}
The Weil pairing is a bilinear, alternating pairing from elliptic-curve torsion points to roots of unity. \par\noindent\textbf{Applications:} it reveals elliptic-curve structure and supports pairing-based cryptography.

\term{Weil Transform}
The Weil transform is an integral transform used to relate functions or representations in harmonic analysis. \par\noindent\textbf{Applications:} it appears in automorphic forms, representation theory, and number theory.

\term{Weierstrass Function}
The Weierstrass function is a function continuous everywhere but differentiable nowhere. \par\noindent\textbf{Applications:} it demonstrates the difference between continuity and differentiability and provides a key example in real analysis.

\term{Weierstrass Product Theorem}
The Weierstrass product theorem represents an entire function through its zeros and an exponential factor, with convergence factors when needed. \par\noindent\textbf{Applications:} it constructs entire functions and studies their zeros in complex analysis.

\term{Weierstrass Theorem (Approximation)}
The Weierstrass approximation theorem states that every continuous function on a closed interval can be uniformly approximated by polynomials. \par\noindent\textbf{Applications:} it justifies polynomial approximation in numerical analysis, modeling, and computation.

\term{Weight (Representation Theory)}
A weight in representation theory is a linear functional describing how a Cartan subalgebra acts on a vector. The collection of weights helps characterize a representation. \par\noindent\textbf{Applications:} weights classify symmetries in Lie theory and quantum mechanics.

\term{Weight (Graph Theory)}
The weight of an edge is its assigned numerical value; the weight of a subgraph is often the sum of its edge weights. \par\noindent\textbf{Applications:} weights represent distance, cost, capacity, or reliability in network optimization.

\term{Well-Defined}
A definition is well-defined when its result does not depend on choices made during the construction. \par\noindent\textbf{Applications:} well-definedness is essential for quotient objects, coordinate definitions, and mathematical functions.

\term{Well-founded Relation}
A relation is well-founded when every nonempty subset has a minimal element under the relation. \par\noindent\textbf{Applications:} well-foundedness guarantees termination of recursive definitions and supports induction arguments.

\term{Well-Ordered}
A well-order is a total order in which every nonempty subset has a least element. \par\noindent\textbf{Applications:} well-orders support transfinite induction and constructions in set theory.

\term{Well-ordering Principle}
The well-ordering principle states that every set can be equipped with a well-order; for positive integers this is the least-number principle. \par\noindent\textbf{Applications:} it supports induction, ordinal methods, and the axiom-of-choice framework.

\term{Well-Posed Problem}
A problem is well-posed when a solution exists, is unique, and changes continuously with the data. \par\noindent\textbf{Applications:} well-posedness is a basic requirement for reliable models and numerical solutions of PDEs and inverse problems.

\term{Weyl Group}
A Weyl group is generated by reflections associated with the roots of a root system. It describes the root system's symmetries. \par\noindent\textbf{Applications:} Weyl groups classify Lie algebras, reflection groups, and representations.

\term{Weyl Algebra}
A Weyl algebra is generated by polynomial variables and their formal derivatives with the relation $[\partial,x]=1$. \par\noindent\textbf{Applications:} it describes differential operators and supports D-modules, PDEs, and representation theory.

\term{Weyl Character Formula}
The Weyl character formula gives the character of an irreducible representation of a semisimple Lie algebra using its highest weight and Weyl group. \par\noindent\textbf{Applications:} it computes representation dimensions and symmetry data.

\term{Weyl Tensor}
The Weyl tensor is the trace-free part of the Riemann curvature tensor. It describes gravitational distortion not determined by local matter density. \par\noindent\textbf{Applications:} it models tidal forces and gravitational waves in relativity.

\term{Weyl's Law}
Weyl's law gives the leading asymptotic growth of Laplacian eigenvalues in terms of the volume of a domain. \par\noindent\textbf{Applications:} it connects spectral data with geometry and is used in PDEs and mathematical physics.

\term{Weyl's Inequality}
Weyl's inequalities bound the eigenvalues of a sum of Hermitian matrices using the eigenvalues of the summands. \par\noindent\textbf{Applications:} they support matrix perturbation theory, optimization, and quantum mechanics.

\term{Whitney Embedding Theorem}
The Whitney embedding theorem states that every smooth $n$-manifold can be smoothly embedded in $\mathbb R^{2n}$. \par\noindent\textbf{Applications:} it represents abstract manifolds as concrete Euclidean subspaces for geometry and topology.

\term{Wigner D-matrix}
Wigner D-matrices represent three-dimensional rotations in bases of angular-momentum states. \par\noindent\textbf{Applications:} they are used in quantum mechanics, spherical harmonics, molecular physics, and 3D orientation.

\term{Wigner's Semicircle Law}
Wigner's semicircle law states that the normalized eigenvalue distribution of a large random symmetric matrix approaches a semicircle. \par\noindent\textbf{Applications:} it supports random matrix theory, statistics, wireless systems, and mathematical physics.

\term{Winding Number}
The winding number is the integer counting how many net times a closed curve winds around a point. \par\noindent\textbf{Applications:} it detects zeros and poles, classifies loops, and appears in complex analysis and topology.

\term{Witt Vector}
Witt vectors package sequences of elements of a ring of characteristic $p$ into a ring of characteristic zero. \par\noindent\textbf{Applications:} they are used in $p$-adic arithmetic, algebraic geometry, and cohomology.

\term{Witt Group}
The Witt group classifies quadratic forms up to adding hyperbolic planes. \par\noindent\textbf{Applications:} it studies quadratic forms, algebraic invariants, and arithmetic geometry.

\term{Word}
A word is a finite sequence of symbols from an alphabet. \par\noindent\textbf{Applications:} words are the basic objects of formal languages, string algorithms, coding, and automata theory.

\term{Wronskian}
The Wronskian is the determinant formed from functions and their successive derivatives. A nonzero Wronskian at a point proves linear independence. \par\noindent\textbf{Applications:} it tests solutions of linear ODEs and helps construct fundamental solution sets.

\term{Well-Ordering Theorem}
The well-ordering theorem states that every set admits a well-order, and it is equivalent to the axiom of choice and Zorn's lemma. \par\noindent\textbf{Applications:} it supports transfinite induction and existence proofs in set theory and algebra.
